\documentclass[10pt,a4paper]{amsart}
\usepackage[margin=19mm]{geometry}
\usepackage{amsmath,amssymb,mathtools}
\usepackage[scr=rsfs,cal=euler]{mathalfa}
\usepackage{xfrac}
\usepackage[unicode,hidelinks,bookmarks=false]{hyperref}
\newcommand{\ie}{i.e.\ }
\newcommand{\cf}{cf.\ }
\newcommand{\resp}{respectively\ }
\newcommand{\Subsection}{\subsection}
\providecommand{\nobreakdash}{\nobreak\hbox{-}}
\setlength{\emergencystretch}{2em}
\multlinegap0pt
\usepackage{mathtools}
\usepackage{bm}
\usepackage{float}
\usepackage{xcolor}
\usepackage{amssymb}
\usepackage[shortlabels]{enumitem}
\usepackage{tikz}
\usepackage[normalem]{ulem}
\usepackage{xfrac}
\newtheorem{lemma}{Lemma}
\newtheorem{proposition}{Proposition}
\newcommand{\Fq}{\mathbb{F}_{q}}
\newcommand{\Fqk}{\mathbb{F}_{q^k}}
\newcommand{\Fqn}{\mathbb{F}_{q^n}}
\DeclareMathOperator{\Gr}{Gr}
\newcommand{\calB}{\mathcal{B}}
\newcommand{\calV}{\mathcal{V}}
\newcommand{\gbinom}[2]{\genfrac{[}{]}{0pt}{}{#1}{#2}}
\title{Decoding Desarguesian spread codes beyond \\half minimum distance}
\author{Ermes Franch, Chunlei Li, Angelica Piccirillo}
\date{Source: arXiv:2607.16890v1 (CC BY 4.0)}
\begin{document}
\maketitle

\noindent\emph{Source and licence.} Decoding Desarguesian spread codes beyond \\half minimum distance. Version arXiv:2607.16890v1 (CC BY 4.0), distributed by arXiv under the Creative Commons Attribution 4.0 International licence, http://creativecommons.org/licenses/by/4.0/, as recorded in the arXiv licence metadata for this exact version and shown on the abstract page https://arxiv.org/abs/2607.16890v1. Full licence notice: \texttt{licenses/arxiv-2607.16890-CC-BY-4.0.txt}.\par\smallskip

\noindent\textit{Editorial scope.} Counted unit: the article's proposition bound (source0.tex lines 1564--1571). The article's complete original proof of it is delivered with it (source0.tex lines 1572--1620, one \texttt{proof} environment of 49 lines). No mathematics is written, repaired or re-derived: the statement and the proof are the article's own lines, in the article's own notation, and no retained line is substituted or re-flowed.  Retained uncounted as the source-local context the proof needs: lines 1543--1549.  Nothing was omitted from any retained span, so the package carries no omission record.  No retained line contains a citation, so the wrapper carries no bibliography and no bibliography entry.  No retained line contains a figure environment, an image-inclusion command or drawing code, so no image is delivered and the package declares no asset dependency.  Editorial formatting only, in addition: an AMS \texttt{amsart} preamble that restates the article's own declarations and macros used by the retained lines, namely the environments \texttt{lemma}, \texttt{proposition} on separate counters, together with the standard packages those lines need.  The retained environments print in this document as Lemma 1, Proposition 1 in order of appearance, whereas the article prints the same items with the numbering of its own full document.  The complete native source of the article as distributed in the arXiv e-print for this exact version is bundled unchanged as \texttt{sources/205511/source0.tex}.
\begin{lemma}
        \label{lemma: counting intersections}
        Let $d_1,d_2,h,n$ be positive integers such that $0\leq h\leq \min\{d_1,d_2\}$ and $d_2-h\leq n-d_1$. Fix a $d_1$-dimensional subspace $\calV_1\leq\Fqn$. Then 
        \[
        \left\lvert\{d_2\text{-dimensional subspaces }\calV_2\leq\Fqn\text{ with }\mathrm{dim}(\calV_1\cap\calV_2)=h\}\right\rvert=q^{(d_1-h)(d_2-h)}\gbinom{n-d_1}{d_2-h}_q\gbinom{d_1}{h}_q.
        \]
    \end{lemma} 

    Thanks to the latter we can provide the following proposition. \begin{proposition}
        \label{ref: our lower bound}
        Let $0\leq i\leq\left\lfloor\frac{\eta kr}{\eta+1}\right\rfloor$, then
        \[
            \left\lvert S_{i,\eta}\right\rvert\geq \gbinom{n}{i}_q-\frac{q^n-1}{q^k-1}\left(\gbinom{n}{i}_q-\sum_{h=0}^\eta q^{(k-h)(i-h)}\gbinom{n-k}{i-h}_q\gbinom{k}{h}_q\right)\eqqcolon \mathrm{LB}_{\eta}.
        \]
        This bound is tight if and only if $i \leq 2 \eta + 1$.
    \end{proposition}

    \begin{proof}
        Let $a\in\Fqn^*$, $h\in\{0,\ldots,\eta\}$ and define 
    \[
    \Sigma_{a,h}\coloneqq \{\calB \in  \Gr_{\Fq}(i,\Fqn) \mid \dim(\calB \cap a \Fqk)=h\}, 
    \]
    then
    \[
    S_{i,\eta}=\bigcap_{a\in\Fqn^*}\left(\bigcup_{h=0}^\eta\Sigma_{a,h}\right).
    \]
    Fix $\calB \in  \Gr_{\Fq}(i,\Fqn)$, the latter follows from the fact that for all $a \in \Fqn^*,~\dim(\calB \cap a \Fqk) \leq \eta$ if and only if for all $a \in \Fqn^*$, there exists $h\in\{0,\ldots,\eta\}$ such that $\dim(\calB \cap a \Fqk)=h$. We first note that distinct values of $a$ can correspond to the same element of the spread. Thus, we define $\mathfrak{A}\coloneqq\Fqn^*/\sim_{q^k}$ where $\sim_{q^k}$ denotes the following equivalence relation on $\Fqn^*$: 
    \[
    a\sim_{q^k} b\text{ if and only if }a\Fqk=b\Fqk.
    \]
    Therefore, $S_{i,\eta}$ can be also seen as
    \[
    S_{i,\eta}=\bigcap_{a\in\mathfrak{A}}\left(\bigcup_{h=0}^\eta\Sigma_{a,h}\right).
    \]
    Recall we want to determine $|S_{i,\eta}|$. Once fixed $a\in\mathfrak{A}$, since for all $h,j\in\{0,\ldots,\eta\}$ with $h\neq j$, $\Sigma_{a,h}\cap\Sigma_{a,j}=\emptyset$, and thanks to Lemma \ref{lemma: counting intersections}, we have that
    \[
    \left\lvert\bigcup_{h=0}^\eta\Sigma_{a,h}\right\rvert=\sum_{h=0}^\eta\left\lvert\Sigma_{a,h}\right\rvert=\sum_{h=0}^\eta q^{(k-h)(i-h)}\gbinom{n-k}{i-h}_q\gbinom{k}{h}_q.
    \]
    As dealing with the intersections for $a$ varying in $\mathfrak{A}$ is nontrivial, we consider the complement with respect to the set $\Gr_{\Fq}(i,\Fqn)$
    \[
    S_{i,\eta}^c=\bigcup_{a\in\mathfrak{A}}\left(\bigcup_{h=0}^\eta\Sigma_{a,h}\right)^c.
    \]
    Therefore 
    \begin{equation}
    \label{eq: evasive space inequality bound}
        \begin{aligned}
            \left\lvert S_{i,\eta}^c\right\rvert&=\left\lvert \bigcup_{a\in\mathfrak{A}}\left(\bigcup_{h=0}^\eta\Sigma_{a,h}\right)^c\right\rvert\leq \sum_{a\in\mathfrak{A}} \left\lvert \left(\bigcup_{h=0}^\eta\Sigma_{a,h}\right)^c\right\rvert\\
            &=\frac{q^n-1}{q^k-1}\left(\gbinom{n}{i}_q-\sum_{h=0}^\eta q^{(k-h)(i-h)}\gbinom{n-k}{i-h}_q\gbinom{k}{h}_q\right),
        \end{aligned}
    \end{equation}
    where $\frac{q^n - 1}{q^k - 1} = |\mathfrak{A}|.$
    From the relation $|S_{i,\eta}| = \gbinom{n}{i}_q -  |S_{i,\eta}^c| $ we obtain the desired result.
    \par In (\ref{eq: evasive space inequality bound}) equality is satisfied if and only if all sets $\left(\bigcup_{h=0}^\eta\Sigma_{a,h}\right)^c$ are disjoint for $a \in \mathfrak{A}$.
    A necessary and sufficient condition for this to be true is given by $i \leq 2 \eta + 1$.
    We can describe the set $\left(\bigcup_{h=0}^\eta\Sigma_{a,h}\right)^c$ as the set of all subspaces of dimension $i$ that intersect $a \Fqk$ with dimension at least $\eta + 1$.
    For $i = 2 \eta + 2$ we can easily choose two subspaces of dimension $\eta + 1$ from two distinct spaces $a \Fqk \neq b \Fqk,$ their sum will have dimension $2 \eta + 2$ and will lie in the intersection of $\left(\bigcup_{h=0}^\eta\Sigma_{a,h}\right)^c \cap \left(\bigcup_{h=0}^\eta\Sigma_{b,h}\right)^c$ showing the condition is necessary.
    
    To show it is also sufficient, consider $\calB \in \left(\bigcup_{h=0}^\eta\Sigma_{a,h}\right)^c$ of dimension $\dim(\calB) = i \leq 2 \eta + 1$ and a subspace of the form $b \Fqk \neq a \Fqk$.
    Since $a \Fqk \cap b \Fqk = \{0\}$ then $\calB \cap a \Fqk$ and $\calB \cap b \Fqk$ will have zero intersection, while their sum will have, at most, dimension $i \leq 2\eta + 1$.
    From which
    \[
      \dim(\calB \cap b \Fqk) \leq  2 \eta + 1  - \dim(\calB \cap  a\Fqk) \leq \eta, 
    \]
    where for the last inequality we used $\dim(\calB \cap a \Fqk) \geq \eta + 1.$
    This ensures that $\calB \notin \left(\bigcup_{h=0}^\eta\Sigma_{b,h}\right)^c$ for any $b \neq a \in \mathfrak{A}$, implying that sets $\left(\bigcup_{h=0}^\eta\Sigma_{b,h}\right)^c$ are always disjoint for $i \leq 2 \eta +1$.
    \end{proof}

\end{document}
